% Modules syllabus à l'échelle d'une école : temps des pages, stress, ruée, traitements lourds,
% et audit des parcours des enseignants, de la DEVE et des étudiants.
% Compilation : latexmk -pdf test-syllabus-charge.tex
\documentclass[11pt,a4paper]{article}
\usepackage[french]{babel}
\usepackage[T1]{fontenc}
\usepackage[utf8]{inputenc}
\usepackage{lmodern,microtype,booktabs,tabularx,longtable,xcolor,hyperref,array}
\usepackage[margin=2.2cm]{geometry}
\hypersetup{colorlinks=true, linkcolor=blue!50!black, urlcolor=blue!50!black}
\newcommand{\code}[1]{\texttt{\detokenize{#1}}}
\newcolumntype{R}{>{\raggedleft\arraybackslash}X}
\title{Modules syllabus à l'échelle d'une école :\\ charge, parcours et ergonomie}
\author{Sébastien Lefebvre}
\date{Novembre 2026 -- version 1.4.0 des modules}

\begin{document}
\maketitle

\section{Objet et méthode}

Les modules \code{mod_syllabus} (fiche et évaluation de chaque module d'enseignement),
\code{mod_syllabusstaff} (suivi des enseignants et de la DEVE) et \code{mod_syllabusstudent}
(« Mes modules » des étudiants) ont été soumis aux mêmes épreuves qu'EP Manager : temps de
chaque page, montée en charge jusqu'à saturation, endurance, ruée, traitements lourds
simultanés, tests unitaires, et validation de bout en bout des parcours. Les parcours ont en
outre été audités pour leur simplicité : nombre de pages et de clics, champs à remplir, clarté
des pages.

\paragraph{Serveur.} Celui des autres tests : 4 c\oe urs, 15~Go, nginx, PHP-FPM~8.3 (16
processus, OPcache), PostgreSQL~16, Moodle~4.5. Le générateur de charge tourne sur la même
machine.

\paragraph{Données} (\code{syllabus_charge_setup.php}, 4 minutes). Une école complète :
\begin{itemize}
  \item cinq années d'étude de douze modules, six par semestre (S1 à S10), soit 60 cours de
    module, chacun avec son activité « Syllabus du module » et deux responsables (120 enseignants) ;
  \item 1\,000 étudiants, 200 par année, inscrits aux douze modules de leur année et à leur cours
    de promotion (une activité « Mes modules » par promotion) ;
  \item la DEVE gestionnaire de la catégorie du cursus ; un cours d'équipe avec le suivi des
    modules ;
  \item des fiches publiées (8 rubriques, 8 séances), deux questions d'enseignant par module ;
  \item un référentiel de 6 domaines et 48 compétences, 6 compétences par module avec un niveau
    croissant au fil des années, un stage sans module par année ;
  \item les évaluations du premier semestre de chaque année closes avec leurs réponses (4\,529
    réponses, 75~\% de participation).
\end{itemize}

\section{Temps des pages}

Chaque page, ouverte trois fois après une première ouverture (\code{syllabus_pages.py}) ; durée
médiane et lectures en base, avant et après les corrections de la section~\ref{sec:parcours}.

\begin{center}
\small
\begin{tabularx}{\linewidth}{@{}l X r r r r@{}}
\toprule
& & \multicolumn{2}{c}{\textbf{Durée}} & \multicolumn{2}{c}{\textbf{Lectures}} \\
\cmidrule(lr){3-4}\cmidrule(l){5-6}
\textbf{Rôle} & \textbf{Page} & \textbf{avant} & \textbf{après} & \textbf{avant} & \textbf{après} \\
\midrule
Étudiant & Mes modules (12 modules) & 0,10 s & 0,09 s & 66 & 66 \\
Étudiant & Fiche du module ; fiche en PDF & 0,10 s ; 0,15 s & 0,08 s ; 0,11 s & 61 & 61 \\
Enseignant & Mes modules (suivi) & 0,07 s & 0,08 s & 52 & 52 \\
Enseignant & Page du module (avec « Où en est ce module ») & 0,11 s & 0,08 s & 67 & 71 \\
Enseignant & Rédaction de la fiche & 0,10 s & 0,11 s & 60 & 66 \\
Enseignant & Évaluation du module ; résultats & 0,10 s ; 0,09 s & 0,09 s ; 0,09 s & 56 ; 62 & 58 ; 64 \\
Enseignant & Matrice des compétences (54 lignes, 65 cours) & 0,08 s & 0,08 s & 47 & \textit{voir note} \\
Enseignant & Progression d'une compétence & 0,14 s & 0,07 s & 151 & 47 \\
DEVE & Tous les modules (50 par page) & 0,18 s & 0,21 s & 174 & 174 \\
DEVE & Résultats comparés (30 modules) & 0,16 s & 0,20 s & 169 & 169 \\
DEVE & Fiches à relire ; référentiel commun & 0,07 s ; 0,08 s & 0,07 s ; 0,09 s & 40 ; 43 & 40 ; 43 \\
DEVE & Référentiel de compétences (arbre de 54 lignes) & 0,09 s & 0,10 s & 61 & 61 \\
DEVE & Progression d'un domaine & 0,13 s & 0,09 s & 155 & 51 \\
DEVE & Livret : page ; PDF du cursus (60 fiches) & 0,09 s ; 2,5 s & 0,09 s ; 2,2 s & 103 & 44 \\
DEVE & Livret du cursus conservé & 0,05 s & 0,04 s & -- & -- \\
\bottomrule
\end{tabularx}
\end{center}

Toutes les pages répondent en moins de 0,21~s pour un seul utilisateur. La progression et le
livret lisaient la fiche publiée et les responsables module par module : ces lectures sont
désormais groupées (\code{syllabus_get_published_bulk}), d'où trois fois moins de requêtes.
La couverture des domaines dans la matrice, ajoutée avec les corrections, est calculée à partir
de l'arbre déjà chargé, sans requête de plus. « Tous les modules » et les résultats comparés
comptent, pour chaque module, les étudiants appelés à répondre (droit Moodle par cours) : deux
requêtes par ligne, 0,2~s pour 50 lignes, gardées telles quelles car exactes.

\section{Montée en charge}

Des utilisateurs virtuels (\code{syllabus_stress.py}) -- 75~\% d'étudiants, 20~\%
d'enseignants, 5~\% de la DEVE -- ouvrent leurs pages en boucle sans temps de réflexion ;
les étudiants évaluent aussi leurs modules du second semestre (écritures sous le verrou du
module). Paliers de 60~s.

\begin{center}
\begin{tabularx}{\linewidth}{@{}r R R R R R r@{}}
\toprule
\textbf{Utilisateurs} & \textbf{pages/s} & \textbf{médiane} & \textbf{95\,\%} & \textbf{max.} &
  \textbf{évaluation} & \textbf{erreurs} \\
\midrule
25  & 48,9 & 0,49 s & 0,70 s & 1,1 s  & 0,72 s & 0 \\
50  & 50,5 & 0,93 s & 1,37 s & 1,9 s  & 1,57 s & 0 \\
100 & 48,0 & 1,91 s & 3,40 s & 4,1 s  & 3,53 s & 0 \\
200 & 49,1 & 3,74 s & 7,12 s & 7,8 s  & 7,24 s & 0 \\
300 & 49,5 & 5,52 s & 10,8 s & 13,3 s & 10,9 s & 0 \\
400 & 55,4 & 6,84 s & 13,4 s & 14,6 s & 13,4 s & 0 \\
\bottomrule
\end{tabularx}
\end{center}

Le processeur est saturé dès 25 utilisateurs virtuels : le débit plafonne à une cinquantaine de
pages par seconde, comme pour EP Manager, et le temps d'attente croît ensuite avec le nombre
d'utilisateurs, régulièrement, sans erreur ni délai dépassé jusqu'à 400. Toutes les pages
attendent à peu près autant : aucune ne se distingue sous charge. L'envoi d'une évaluation compte
deux requêtes (l'envoi et la page de retour). 1\,136 évaluations ont été enregistrées pendant le
test, sans réponse orpheline ni participation sans réponse.

ENDURANCE

\section{Ruée à l'ouverture des évaluations}

Les 1\,000 étudiants évaluent \textbf{au même instant} le même module de leur année : 200
étudiants par module, en concurrence sur le verrou du module et sur le figeage des questions,
qui a lieu à la première ouverture du formulaire (\code{syllabus_ruee.py}). Chacun ouvre « Mes
modules », le formulaire, puis l'envoie.

\begin{center}
\begin{tabularx}{\linewidth}{@{}X r@{}}
\toprule
Évaluations confirmées & 1\,000 / 1\,000 \\
Durée jusqu'à la dernière ; médiane et 95\,\% par étudiant & 54 s ; 46 s et 53 s \\
Participations et réponses en base & 1\,000 et 1\,000 \\
Doublons ; réponses obligatoires manquantes ; erreurs & 0 ; 0 ; 0 \\
Modules figés (questions copiées une seule fois) & 5 / 5 \\
\bottomrule
\end{tabularx}
\end{center}

Les 3\,000 requêtes sont servies au débit maximal du serveur ; aucune n'échoue et l'anonymat
comme l'unicité des réponses tiennent. Une ouverture réelle s'étale sur des heures : 200
étudiants qui répondent dans la même minute attendent quelques secondes.

\section{Traitements lourds simultanés}

Pendant que 50 utilisateurs virtuels ouvrent leurs pages, sont lancés ensemble
(\code{syllabus_lourds.py}) :

\begin{center}
\begin{tabularx}{\linewidth}{@{}X r r@{}}
\toprule
\textbf{Traitement} & \textbf{Durée} & \textbf{Résultat} \\
\midrule
Tâche d'ouverture des évaluations du second semestre (30 modules) & 65 s & 4\,825 messages \\
Livret de tout le cursus, par la DEVE (60 fiches) & 13,0 s & PDF de 778 Ko \\
Livret de chaque année, par cinq enseignants & 3,6 à 4,7 s & PDF de 300 Ko \\
Exports CSV : matrice, résultats comparés, progression d'un domaine & 1,0 à 1,3 s & conformes \\
\midrule
Charge interactive pendant ce temps (médiane ; 95\,\% ; max.) & \multicolumn{2}{r}{1,00 s ; 1,47 s ; 2,9 s} \\
Même charge sans traitement lourd & \multicolumn{2}{r}{0,93 s ; 1,37 s ; 1,9 s} \\
\bottomrule
\end{tabularx}
\end{center}

Les pages des autres utilisateurs n'en sont presque pas ralenties. La tâche d'ouverture envoie
environ 75 messages par seconde (9~ms par message hors charge) : elle tourne tous les quarts
d'heure, en tâche de fond. Le livret composé est conservé : le suivant est servi en 0,04~s.

\paragraph{Mesure corrigée.} Un premier passage donnait 11 à 13~s pour les exports CSV : le
script les lançait tous dans la même session de la DEVE, et Moodle sert l'une après l'autre les
requêtes d'une même session (verrou de session). Avec une session par traitement, ils prennent
1 à 1,3~s.

PROFIL

\section{Parcours et ergonomie}\label{sec:parcours}

\paragraph{Méthode.} Une revue indépendante, conduite sur le code, les chaînes et les captures
réelles, a reconstitué chaque parcours (enseignant, DEVE, étudiant), compté pages, clics et
champs, et relevé 60 constats classés par priorité. Chaque constat a été vérifié dans le code ;
les corrections retenues ont ensuite été validées par le scénario de bout en bout, rejoué par
les pages web (section~\ref{sec:validation}).

\subsection{Gains sur les parcours}

\begin{center}
\small
\begin{tabularx}{\linewidth}{@{}X l l@{}}
\toprule
\textbf{Parcours} & \textbf{Avant} & \textbf{Après} \\
\midrule
Étudiant : évaluer ses 12 modules depuis « Mes modules » & 46 clics, 46 pages & 24 clics, 24 pages \\
Étudiant : formulaire d'évaluation (5 questions) & HAUTAVANT & HAUTAPRES \\
Enseignant : rédiger la fiche avec 3 enregistrements intermédiaires & 8 clics, 8 pages & 5 clics, 5 pages \\
Enseignant : savoir ce qu'il reste à faire sur son module & 4 pages & 1 encadré, liens directs \\
DEVE : relire 12 fiches depuis « Fiches à relire » & 57 clics & 24 clics \\
DEVE : saisir 8 sous-compétences d'un domaine & 16 champs, 8 choix du parent & 8 champs \\
DEVE : matrice de toutes les années & filtre trompeur (« Non précisée ») & « Toutes les années » \\
\bottomrule
\end{tabularx}
\end{center}

\subsection{Corrections}

\paragraph{Étudiants.}
\begin{itemize}
  \item Formulaire d'évaluation : échelles et oui/non sur une ligne (au lieu de cinq lignes par
    question) ; rappel « vous ne pourrez répondre qu'une fois » près du bouton d'envoi ; message
    d'erreur précis (« questions marquées d'un astérisque »).
  \item Après l'envoi (ou l'annulation), retour à « Mes modules » au lieu de la fiche : le module
    suivant est à un clic.
  \item « Mes modules » range les modules par année et période, avec un intertitre par période ;
    le nom du cours n'est plus répété quand il est celui du module.
  \item L'état d'une évaluation passée s'appelle « Terminée » (et non « Close », qui se lisait
    comme l'anglais).
\end{itemize}

\paragraph{Enseignants responsables.}
\begin{itemize}
  \item La page du module s'ouvre sur « Où en est ce module » : fiche (publiée, en cours,
    renvoyée avec le commentaire de la DEVE), compétences visées, questions de l'enseignant et
    état de l'évaluation, chacune avec son bouton ; les caractéristiques manquantes (crédits,
    période, année) sont signalées avec un lien vers les paramètres.
  \item « Enregistrer le brouillon » reste sur la rédaction ; la soumission exige les rubriques
    essentielles (objectifs, acquis, contenu, modalités d'évaluation) et au moins une séance ou
    des heures de travail personnel : une fiche vide n'arrive plus à la DEVE.
  \item Séances : intitulé « Séance 1, 2\dots », champs avec leur indication (intitulé, heures,
    intervenant extérieur), travail personnel saisi avec les séances ; la langue est préremplie
    (« Français »). La DEVE n'est plus proposée comme intervenant ni comme responsable.
  \item Période : choisie parmi celles déjà utilisées dans l'école, ou saisie ; « S5 » et
    « s5 » ne font plus deux colonnes de la frise ou du livret.
  \item Compétences : le formulaire n'est plus précédé du même tableau ; un domaine se vise par
    ses sous-compétences ; bouton « Enregistrer » aussi en tête de l'arbre.
  \item Évaluation : le message des questions figées suit l'état réel (il annonçait une
    évaluation « ouverte » quand elle était terminée) ; lien vers les résultats ; « Modifier »
    amène au formulaire ; suppression d'une question confirmée.
  \item Résultats : taux de réponse en pourcentage et tableau des moyennes avant le détail.
  \item Navigation du module dans l'ordre du travail : rédiger, compétences, évaluation,
    résultats, responsables. Titres d'onglet du navigateur distincts pour chaque page.
  \item Suivi : le nom du module ouvre sa fiche ; les modules qui attendent le responsable sont
    en tête et surlignés ; le lien « Résultats » n'apparaît que s'ils sont visibles.
\end{itemize}

\paragraph{DEVE.}
\begin{itemize}
  \item Relecture : décision sous la fiche relue, rubriques et séances modifiées depuis la
    version publiée signalées, aide sur Publier / Renvoyer / Refuser, commentaire exigé par le
    navigateur pour renvoyer ou refuser ; après la décision, retour à « Fiches à relire ». La
    liste indique responsables et « Première version » ou « Mise à jour », et signale les cours
    où la DEVE n'a pas le droit de décider.
  \item Fenêtre en lot : « Appliquer » exige une ouverture ; un bouton distinct, confirmé,
    retire l'évaluation ; les cases des cours sans le rôle de la DEVE sont grisées ; une
    sélection vide est refusée.
  \item Référentiel commun : « Enregistrer les réglages » ; « Monter » et « Descendre » seulement
    s'il y a une voisine ; suppressions confirmées.
  \item Référentiels de compétences : après un ajout, le formulaire reste ouvert sous le même
    parent, avec le code suivant proposé (C1.4) ; import et réglages du référentiel repliés ;
    suppressions confirmées.
  \item Matrice : « Matrice des compétences », option « Toutes les années » ; un domaine compte
    les cours qui visent l'une de ses sous-compétences.
  \item Résultats comparés : le nom du module ouvre ses résultats. Livret : « Périmètre du
    livret » ; téléchargement impossible sans fiche publiée.
\end{itemize}

\subsection{Constats non retenus ou reportés}

\begin{itemize}
  \item \emph{« Évaluation » désigne l'enquête auprès des étudiants et les « modalités
    d'évaluation » de la fiche.} Le vocabulaire est celui demandé ; la question est posée au
    responsable du projet (voir la conclusion).
  \item \emph{Fenêtre d'évaluation réglable dans les paramètres et dans « Évaluation du
    module ».} Gardée aux deux endroits, une note le dit dans les paramètres.
  \item \emph{Enregistrer qui a fixé la fenêtre (enseignant ou DEVE)}, \emph{confirmer qu'un
    responsable se retire lui-même}, \emph{colonnes des résultats comparés groupées par type},
    \emph{en-têtes de colonnes de la matrice sur deux lignes} : utiles mais secondaires, reportés.
  \item \emph{Crédits obligatoires dans les paramètres} : un module peut légitimement ne pas en
    porter ; ils sont signalés dans « Où en est ce module » au lieu d'être exigés.
\end{itemize}

\section{Tests et validation}\label{sec:validation}

\begin{itemize}
  \item \textbf{Tests unitaires} : 44 (37 pour \code{mod_syllabus}, 6 pour le suivi, 1 pour
    « Mes modules »), dont 12 nouveaux : réglages communs et leurs bornes, types (suppression
    refusée s'ils servent), questions communes et de l'enseignant (contrôles, portée par type,
    ordre, module figé intouchable), état de l'évaluation, messages (DEVE à la soumission,
    responsables à la décision, étudiants n'ayant pas répondu à l'ouverture), événement
    d'évaluation sans réponse, responsables possibles (sans la DEVE), suppression de l'activité
    et d'un cours (rien ne reste en base), périmètre du suivi et onglets selon le rôle, fiche
    incomplète refusée à la soumission, fiches publiées lues en groupe. Test de conformité RGPD
    de Moodle réussi pour les trois modules ; contrôle de style Moodle sans remarque.
  \item \textbf{Scénario de bout en bout} (\code{syllabus_interactions.py}, rapport
    \emph{Test des modules syllabus}) rejoué sur la version 1.4.0 : \textbf{44 contrôles sur 44},
    dont les nouveaux parcours -- retour à « Mes modules » après chaque évaluation (2\,145
    évaluations), décision de la DEVE sous la fiche et retour à la liste des fiches à relire.
\end{itemize}

\paragraph{Défauts du banc trouvés.} La remise à zéro du scénario effaçait tous les référentiels
de compétences du site, y compris celui du jeu de données à l'échelle : elle n'efface plus que
le sien. Le scénario reconnaissait les liens « Évaluer » à leur forme exacte et n'a pas vu ceux
qui portent désormais le retour à « Mes modules » : il les lit maintenant en entier.

\section{Conclusion}

À l'échelle d'une école (60 modules, 1\,000 étudiants, 120 responsables), toutes les pages
répondent en moins de 0,21~s ; le serveur de test plafonne à une cinquantaine de pages par
seconde sans erreur jusqu'à 400 utilisateurs sans temps de réflexion ; 1\,000 évaluations
simultanées s'enregistrent toutes, une seule fois chacune ; les traitements lourds ne gênent pas
les autres utilisateurs. Les parcours ont été raccourcis là où ils faisaient revenir l'utilisateur
en arrière (évaluer, rédiger, relire), et chaque page dit maintenant ce qu'il reste à faire.

Une question reste ouverte : faut-il nommer autrement l'enquête auprès des étudiants
(« Évaluation du module ») pour la distinguer des « Modalités d'évaluation » de la fiche, par
exemple « Avis des étudiants » ?

\section{Reproduire}

{\small
\begin{verbatim}
cd docs/audit/banc
# Scénario de bout en bout d'abord (il recrée le référentiel commun), puis le jeu à l'échelle.
sudo -u www-data php syllabus_charge_setup.php site-x.json > syllabus-charge.json
python3 syllabus_pages.py site-x.json syllabus-charge.json > syllabus-pages.json
python3 syllabus_stress.py paliers site-x.json syllabus-charge.json > syllabus-stress.json
python3 syllabus_stress.py endurance site-x.json syllabus-charge.json 100 30 > endurance.json
python3 syllabus_ruee.py site-x.json syllabus-charge.json > ruee.json
python3 syllabus_lourds.py site-x.json syllabus-charge.json > lourds.json
python3 syllabus_profil.py /home/user/moodle-stress/prof site-x.json syllabus-charge.json
sudo -u www-data php syllabus_charge_setup.php site-x.json reset      # pour rejouer
\end{verbatim}
}

\end{document}
